\paragraph{Problem statement.}
Laser ultrasound (LUS) enables non-contact inspection of critical infrastructure; however, acquired signals are invariably corrupted by intense ambient and photonic noise, yielding signal-to-noise ratios (SNR) as low as $-10$~dB \cite{laser_ultrasound_noise}. This low SNR fundamentally limits defect detection reliability in non-destructive testing (NDT), manifesting in three concrete forms: arrival time jitter that degrades thickness and velocity measurements \cite{arrival_time_jitter}, amplitude distortion that masks fatigue cracks and corrosion \cite{amplitude_distortion}, and elevated false-negative rates that permit dangerous defects to go undetected \cite{false_negative_ndt}. Conventional filtering methods---bandpass, wavelet denoising, and spectral subtraction---operate with fixed priors and cannot adapt to the non-stationary, structured noise characteristic of laser-generated acoustic emissions \cite{conventional_filter_limits}. While deep learning (DL) has demonstrated strong denoising capacity in acoustic imaging tasks \cite{deep_learning_ultrasound}, existing DL architectures trained on additive white Gaussian noise (AWGN) suffer from a critical sim-to-real gap: they over-smooth waveform morphology, blur arrival edges, and suppress legitimate signal components when deployed on structured laboratory or field noise \cite{over_smoothing_issue,sim_to_real_gap}. Preserving waveform fidelity while achieving robust SNR improvement remains an open challenge in LUS-based NDT.
